\documentclass[10pt,a4paper]{amsart}
\usepackage[margin=19mm]{geometry}
\usepackage{amsmath,amssymb,mathtools}
\usepackage[scr=rsfs,cal=euler]{mathalfa}
\usepackage{xfrac}
\usepackage[unicode,hidelinks,bookmarks=false]{hyperref}
\newcommand{\ie}{i.e.\ }
\newcommand{\cf}{cf.\ }
\newcommand{\resp}{respectively\ }
\newcommand{\Subsection}{\subsection}
\providecommand{\nobreakdash}{\nobreak\hbox{-}}
\setlength{\emergencystretch}{2em}
\multlinegap0pt
\usepackage{bm}
\usepackage{bbm}
\usepackage{dsfont}
\usepackage{xspace}
\usepackage{cancel}
\usepackage{mathtools}
\usepackage{amsthm}
\makeatletter
\@ifundefined{proposition}{\newtheorem{proposition}{Proposition}}{}
\makeatother
\DeclareMathOperator{\B}{B}
\DeclareMathOperator{\C}{C}
\DeclareMathOperator{\QAut}{QAut}
\newcommand{\I}{\mathds{1}}
\newcommand{\id}{\mathrm{id}}
\newcommand{\is}[2]{{\left\langle{#1}\,\vline\,#2\right\rangle}}
\newcommand{\ket}[1]{{\left|#1\right\rangle}}
\newcommand{\tens}{\otimes}
\title[Quantum Mycielskians: symmetries, twin vertices and distingu]{Quantum Mycielskians: symmetries, twin vertices and distinguishing labelings}
\author{Arkadiusz Bochniak, Igor Che{\l}stowski, Pawe{\l} Kasprzak, Piotr M. So{\l}tan}
\date{Source: arXiv:2502.10521v2 (CC BY 4.0)}
\begin{document}
\maketitle

\noindent\emph{Source and licence.} Quantum Mycielskians: symmetries, twin vertices and distinguishing labelings. Version arXiv:2502.10521v2 (CC BY 4.0), distributed under the Creative Commons Attribution 4.0 International licence, as recorded in the arXiv licence metadata for this exact version and shown on the abstract page https://arxiv.org/abs/2502.10521v2. Full rights notice: licenses/arxiv-2502.10521-CC-BY-4.0.txt.\par\smallskip

\noindent\textit{Editorial scope.} Counted unit: the Proposition whose source label is \texttt{None} in the article (source0.tex lines 313--315, source label \texttt{None}), delivered together with the article's own complete original proof of it (source0.tex lines 317--375, one \texttt{proof} environment of 59 lines). No mathematics is written, repaired or re-derived: statement and proof are the article's own text line for line. Required context retained uncounted: prop:deltas (lines 305--307). Every definition, display and label of those lines is retained unchanged. Omitted from this package and not redistributed: 1--304, 308--312, 316--316, 376--1439. Nothing inside a retained excerpt is omitted or altered: every retained body is delivered exactly as the article's lines are written, so no omission receipt of any kind accompanies this package. Citations: the retained excerpts carry no citation command at all, so no bibliography entry of the article is transcribed and no reference is redistributed. Reference adaptation: none was needed and none was made; every reference inside the delivered unit resolves inside it, so no unresolved reference or citation is produced. Printed slips kept literal: no printed word, symbol, hypothesis or number of the counted statement or of its proof was altered. Layout and class: the article's own class is \texttt{amsart} with packages including mathtools, xcolor, geometry, dsfont, xy, graphicx, hyperref, mathalfa, tikz, amssymb, amsthm, mathrsfs, microtype, mdframed; the wrapper is the lane's pinned \texttt{amsart} editorial wrapper at 10pt on A4, which restates the theorem-like environments the retained text needs and adds the article's own macro definitions that the retained text needs (\texttt{\textbackslash B}, \texttt{\textbackslash C}, \texttt{\textbackslash I}, \texttt{\textbackslash QAut}, \texttt{\textbackslash id}, \texttt{\textbackslash is}, \texttt{\textbackslash ket}, \texttt{\textbackslash tens}), with extra wrapper packages (bm, bbm, dsfont, xspace, cancel, mathtools). The delivered numbering is the unit's own. Assets: no retained span contains a figure environment, a graphics-inclusion command, a PDF-page inclusion, a table or a TikZ picture; no image file is copied into this package, the package declares no asset dependency and no image file is redistributed with it. Licence: version arXiv:2502.10521v2 of this article is distributed under the Creative Commons Attribution 4.0 International licence, as recorded in the arXiv licence metadata for this exact version and shown on the abstract page https://arxiv.org/abs/2502.10521v2. A full rights notice is delivered at licenses/arxiv-2502.10521-CC-BY-4.0.txt. Characters: every retained excerpt is pure ASCII, so the delivered engine, XeLaTeX with its default Latin Modern text font, reports no missing character.
\begin{proposition}\label{prop:deltas}
Let $\mathcal{F}$ and $\mathcal{G}$ be quantum graphs. If $\mathcal{F}$ and $\mathcal{G}$ are quantum isomorphic, then $\delta_{\scriptscriptstyle\mathcal{F}}=\delta_{\scriptscriptstyle\mathcal{G}}$.
\end{proposition}

\begin{proposition}%\label{prop:iso}
Let $\mathcal{G}$ and $\mathcal{F}$ be quantum isomorphic. Then $\mu_{r-1}(\mathcal{G})$ and $\mu_{r-1}(\mathcal{F})$ are quantum isomorphic for any $r\geq 1$.
\end{proposition}

\begin{proof}
Let
\[
p=[p_{ij}]_{i,j}\in\B\bigl(L^2(\mathcal{G}),L^2(\mathcal{F})\bigr)\tens\mathcal{O}\bigl(\QAut(\mathcal{F}),\QAut(\mathcal{G})\bigr)
\]
be the unitary matrix defining $\mathcal{O}(\QAut(\mathcal{F}),\QAut(\mathcal{G}))$ as above (remember the latter is non-zero) and let
\[
\rho_{\scriptscriptstyle\mathcal{F},\mathcal{G}}\colon\C(\mathcal{G})\ni e_j\longmapsto\sum\limits_i f_i\tens p_{ij}\in\C(\mathcal{F})\tens\mathcal{O}\bigl(\QAut(\mathcal{F}),\QAut(\mathcal{G})\bigr)
\]
be the corresponding unital $*$-homomorphism intertwining the quantum adjacency matrices. By Proposition~\ref{prop:deltas} we know that $\delta_{\scriptscriptstyle\mathcal{G}}=\delta_{\scriptscriptstyle\mathcal{F}}$ and for the remainder of the proof we will denote this value by $\delta$. The maps introduced above can be represented as
\[
p\colon L^2(\mathcal{G})\tens\mathcal{O}\bigl(\QAut(\mathcal{F}),\QAut(\mathcal{G})\bigr)\ni\xi\tens a\longmapsto\sum\limits_{i,j}\ket{f_i}\is{e_j}{\xi}\tens p_{ij}a\in L^2(\mathcal{F})\tens\mathcal{O}\bigl(\QAut(\mathcal{F}),\QAut(\mathcal{G})\bigr)
\]
and
\[
\rho_{\scriptscriptstyle\mathcal{F},\mathcal{G}}(\xi)=p(\xi\tens \I_{\mathcal{O}(\QAut(\mathcal{F}),\QAut(\mathcal{G}))}).
\]
We define an element $\widehat{p}\in\B(L^2(\mu_{r-1}(\mathcal{G})),L^2(\mu_{r-1}(\mathcal{F})))\tens\mathcal{O}(\QAut(\mathcal{F}),\QAut(\mathcal{G}))$:
\[
\widehat{p}=\varpi_{\scriptscriptstyle\mathcal{F},0}\varpi_{\scriptscriptstyle\mathcal{G},0}^*+\sum\limits_{k=1}^r\varpi_{\scriptscriptstyle\mathcal{F},k} p\varpi_{\scriptscriptstyle\mathcal{G},k}^*,
\]
where $\varpi_{\scriptscriptstyle\mathcal{G},j}=\iota_{\scriptscriptstyle\mathcal{G},j}\tens\I_{\mathcal{O}(\QAut(\mathcal{F}),\QAut(\mathcal{G}))}$ and $\varpi_{\scriptscriptstyle\mathcal{F},j}=\iota_{\scriptscriptstyle\mathcal{F},j}\tens\I_{\mathcal{O}(\QAut(\mathcal{F}),\QAut(\mathcal{G}))}$, with $\iota_{\scriptscriptstyle\mathcal{G},j}$ and $\iota_{\scriptscriptstyle\mathcal{F},j}$ ($j=1,\ldots,r$)denoting the isometric embeddings of $L^2(\mathcal{G})$ and $L^2(\mathcal{F})$ into $L^2(\mu_{r-1}(\mathcal{G}))$ and $L^2(\mu_{r-1}(\mathcal{F}))$, respectively.

First, notice that
\[
\begin{aligned}
\widehat{p}\widehat{p}^*&=\varpi_{\scriptscriptstyle\mathcal{F},0}\varpi_{\scriptscriptstyle\mathcal{F},0}^* +\sum\limits_{k=1}^r\varpi_{\scriptscriptstyle\mathcal{F},k} pp^*\varpi_{\scriptscriptstyle\mathcal{F},k}^* =\sum\limits_{j=0}^r\varpi_{\scriptscriptstyle\mathcal{F},j}\varpi_{\scriptscriptstyle\mathcal{F},j}^*=\id_{L^2(\mu_{r-1}(\mathcal{F}))},\\
\widehat{p}^*\widehat{p}&=\varpi_{\scriptscriptstyle\mathcal{G},0}\varpi_{\scriptscriptstyle\mathcal{G},0}^* +\sum\limits_{k=1}^r\varpi_{\scriptscriptstyle\mathcal{G},k} pp^*\varpi_{\scriptscriptstyle\mathcal{G},k}^* =\sum\limits_{j=0}^r\varpi_{\scriptscriptstyle\mathcal{G},j}\varpi_{\scriptscriptstyle\mathcal{G},j}^*=\id_{L^2(\mu_{r-1}(\mathcal{G}))},
\end{aligned}
\]
so that $\widehat{p}$ is unitary.

Next, we check the intertwining property. First, we compute
\begin{align*}
\widehat{p}A_{\scriptscriptstyle\mu_{r-1}(\mathcal{G})}&=\biggl(\varpi_{\scriptscriptstyle\mathcal{F},0}\varpi_{\scriptscriptstyle\mathcal{G},0}^* +\sum\limits_{k=1}^r\varpi_{\scriptscriptstyle\mathcal{F},k} p\varpi_{\scriptscriptstyle\mathcal{G},k}^*\biggr)\biggl(\delta\varpi_{\scriptscriptstyle\mathcal{G},r}\eta_{\scriptscriptstyle\mathcal{G}}\varpi_{\scriptscriptstyle\mathcal{G},0}^*+\delta\varpi_{\scriptscriptstyle\mathcal{G},0}\eta^*_{\scriptscriptstyle\mathcal{G}}\varpi_{\scriptscriptstyle\mathcal{G},r}^*+\varpi_{\scriptscriptstyle\mathcal{G},1} A_{\scriptscriptstyle\mathcal{G}}\varpi_{\scriptscriptstyle\mathcal{G},1}^*\\
&\qquad\qquad\qquad\qquad\qquad\qquad\qquad\qquad+\sum\limits_{k=1}^{r-1}\bigl(\varpi_{\scriptscriptstyle\mathcal{G},k} A_\mathcal{G}\varpi_{\scriptscriptstyle\mathcal{G},k+1}^* +\varpi_{\scriptscriptstyle\mathcal{G},k+1} A_\mathcal{G}\varpi_{\scriptscriptstyle\mathcal{G},k}^*\bigr)\biggr)\\
&=\delta\varpi_{\scriptscriptstyle\mathcal{F},0}\eta_{\scriptscriptstyle\mathcal{G}}^*\varpi_{\scriptscriptstyle\mathcal{G},r}^* +\delta\varpi_{\scriptscriptstyle\mathcal{F},r}p\eta_{\scriptscriptstyle\mathcal{G}}\varpi_{\scriptscriptstyle\mathcal{G},0}^* +\varpi_{\scriptscriptstyle\mathcal{F},1}pA_{\scriptscriptstyle\mathcal{G}}\varpi_{\scriptscriptstyle\mathcal{G},1}^*\\
&\qquad\qquad\qquad\qquad\qquad\qquad\qquad\qquad+\sum\limits_{j=1}^{r-1}\bigl(\varpi_{\scriptscriptstyle\mathcal{F},j}pA_{\scriptscriptstyle\mathcal{G}}\varpi_{\scriptscriptstyle\mathcal{G},j+1}^*+\varpi_{\scriptscriptstyle\mathcal{F},j+1}pA_{\scriptscriptstyle\mathcal{G}}\varpi_{\scriptscriptstyle\mathcal{G},j}^*\bigr),
\end{align*}
and similarly,
\begin{align*}
A_{\scriptscriptstyle\mu_{r-1}(\mathcal{F})}\widehat{p}=\delta\varpi_{\scriptscriptstyle\mathcal{F},0}\eta_{\scriptscriptstyle\mathcal{F}}^* p\varpi_{\scriptscriptstyle\mathcal{G},r}^* +\delta\varpi_{\scriptscriptstyle\mathcal{F},r}\eta_{\scriptscriptstyle\mathcal{F}}\varpi_{\scriptscriptstyle\mathcal{G},0}^* &+\varpi_{\scriptscriptstyle\mathcal{F},1}A_{\scriptscriptstyle\mathcal{F}}p\varpi_{\scriptscriptstyle\mathcal{G},1}^*\\
&+\sum\limits_{j=1}^{r-1}\bigl(\varpi_{\scriptscriptstyle\mathcal{F},j}A_{\scriptscriptstyle\mathcal{F}}p\varpi_{\scriptscriptstyle\mathcal{G},j+1}^*+\varpi_{\scriptscriptstyle\mathcal{F},j+1}A_{\scriptscriptstyle\mathcal{F}}p\varpi_{\scriptscriptstyle\mathcal{G},j}^*\bigr).
\end{align*}
Since $p\eta_{\scriptscriptstyle\mathcal{G}}=\eta_{\scriptscriptstyle\mathcal{F}}$, we also have $\eta_{\scriptscriptstyle\mathcal{G}}^*=\eta_{\scriptscriptstyle\mathcal{F}}^* p$. This shows that indeed the intertwining property holds. Finally, the map
\[
\rho_{\scriptscriptstyle\mu_{r-1}(\mathcal{F}),\mu_{r-1}(\mathcal{G})}
\left(\left[\begin{smallmatrix}
\lambda\\\xi_1\\[-3pt]\vdots\\[1pt]\xi_r
\end{smallmatrix}\right]\right)
=\widehat{p}
\left(\left[\begin{smallmatrix}
\lambda\\\xi_1\\[-3pt]\vdots\\[1pt]\xi_r
\end{smallmatrix}\right]\tens\I_{\mathcal{O}(\QAut(\mathcal{F}),\QAut(\mathcal{G}))}\right)
\]
is a unital $*$-homomorphism $\C(\mu_{r-1}(\mathcal{G}))\to
\C(\mu_{r-1}\mathcal{F}))\tens\mathcal{O}(\QAut(\mathcal{F}),\QAut(\mathcal{G}))$ 
by construction. Since the algebra $\mathcal{O}(\QAut(\mu_{r-1}(\mathcal{F})),\QAut(\mu_{r-1}(\mathcal{G})))$ is universal for such homomorphisms, it has a non-zero homomorphic image, so it is itself non-zero. It follows that $\mu_{r-1}(\mathcal{G})$ and $\mu_{r-1}(\mathcal{F})$ are quantum isomorphic.
\end{proof}

\end{document}
